\documentclass[conference]{IEEEtran}
\IEEEoverridecommandlockouts
% The preceding line is only needed to identify funding in the first footnote. If that is unneeded, please comment it out.
\usepackage{cite}
\usepackage{amsmath,amssymb,amsfonts}
\usepackage{algorithmic}
\usepackage{graphicx}
\usepackage{textcomp}
\usepackage{xcolor}
\def\BibTeX{{\rm B\kern-.05em{\sc i\kern-.025em b}\kern-.08em
		T\kern-.1667em\lower.7ex\hbox{E}\kern-.125emX}}
\begin{document}
	
	\title{Layer-wise Analysis of Representation Fine-Tuning under resource constraints: Which layers matter most?}
	
	\author{\IEEEauthorblockN{Muhammad Saad Azam}
		\IEEEauthorblockA{\textit{Department of Computer Science} \\
			\textit{University of Engineering and Technology}\\
			Lahore, Pakistan \\
			saad.azam1090@gmail.com}}
	
	\maketitle
	
	\begin{abstract}
		This document is a model and instructions for \LaTeX.
		This and the IEEEtran.cls file define the components of your paper [title, text, heads, etc.]. *CRITICAL: Do Not Use Symbols, Special Characters, Footnotes, 
		or Math in Paper Title or Abstract.
	\end{abstract}
	
	\begin{IEEEkeywords}
		component, formatting, style, styling, insert
	\end{IEEEkeywords}
	
	\section{Introduction}
	
	Since the introduction of transformers \cite{vaswani2017attention}, the landscape has changed to building large-scale language models. Today, Large Language Models (LLMs) are comprised of billions of parameters \cite{brown2020language}. With this enormous size comes the challenging task of fine-tuning these models for specific tasks. While full finetuning still yields more accurate results, the cost of this is enormous. To address this challenge, the field of parameter-efficient finetuning (PEFT) has gathered considerable attention. These methods enable researchers to experiment on state-of-the-art models with limited computational resources.
	
	Different techniques have been introduced to fine-tune LLMs. Adapter layers insert a small module in a frozen transformer and train it \cite{houlsby2019parameter}. Prefix tuning modifies the inputs to attentions instead of modifying weights or activations \cite{li2021prefix}. Hu et al. \cite{hu2022lora} introduced a Low-Rank Adaptation (LoRA) of Large Language Models, a way to fine-tune a model without fine-tuning all of its parameters. All these methods either change the inputs or modify the weights without requiring full finetuning.
	
	One more highly efficient technique of PEFT is Representation Finetuning (ReFT) for Language Models. In ReFT, the representations are directly modified instead of weights of the models \cite{wu2024reft}. Hence, ReFT acts as a drop-in intervention in representations and steer the output of the model towards a specific direction. Furthermore, ReFT trains over less parameters compared to a standard LoRA setup and achieves fairly competent performance over various tasks.
	
	Unlike weight-based adapters which modify model's knowledge base, ReFT operates on the hidden states. ReFT helps the model learn geometric transformations which in result nudge the direction of the model's output without altering its memory. This distinction is crucial because it signals the success of these representations depends more on the semantic richness of the representations at a given depth, which can vary depending on the model architecture and the task of finetuning.
	
	However, ReFT finetuning has a large hyperparameter search space and it is sensitive to layer selection \cite{wu2024reft}. Also, if ReFT is applied to multiple layers, the intervention matrices must be kept active during inference, introducing some inference delay. Furthermore, ReFT has not been thoroughly studied for Small Language Models (SLMs) under resource constraints on a limited dataset. The original paper also showcases how ReFT is susceptible to nature of tasks as well, further expanding the hyperparameter search space. This requires investigation on how the search space can be narrowed down based on the model architecture and the dataset constraints in place.
	
	With this backdrop, it can be highlighted that the efficacy of ReFT is regulated by layer-wise specialization and architectural constraints. Existing work has primarily focused on 7B+ models. The behavior of ReFT in the 'Low-Resource' areas  is under-explored. The research will focus on the following questions:
	\begin{itemize}
		\item What does the trend of loss looks like for the layers? Do some layers tend to overfit in ReFT while other do not? Does the training comply with existing literature on layer fine-tuning?
		\item Does the accuracy change drastically between the layers to specific tasks?
		\item Are interventions productive for small language models?
		\item Are the accuracy and representation trends between the models similar across all layers? 
	\end{itemize}
	
	\section{Literature Review}
	
	\subsection{Representation Finetuning}
	Representation Finetuning (ReFT) is a parameter-efficient finetuning technique for LLMs \cite{wu2024reft}. It trains interventions over the hidden states to steer model behavior into the direction of the finetuning task at inference time. Unlike other methods that target weights, ReFT modifies the hidden represenations. The main inspiration of this methodology is the study of causal interventions to direct the behavior of the model \cite{geiger2023causal}. There are two types of interventions introduced originally: LoReFT and DiReFT. LoReFT learns intervention in linear subspace spanned by a low-rank projection matrix, while DiReFT simply shifts the representation by a fixed vector. DiReFT is more straightforward and efficient but is not as accurate as LoReFT is. LoReFT transformation looks like this: 
	$$\Phi(h) = h + R^T (Wh + b - Rh)$$
	Where $h$ is the original hidden representation, $R$ is the low-rank projection matrix, $W$ is the learned weight matrix and $b$ is the learned bias vector.
	
	\subsection{Linearly Decodable Features}
	Research has shown that many features encoded in the hidden representations are linearly decodable, simulated with linear probing experiments \cite{gurnee2023space}. This is the area of interpretation that ReFT targets. When features are rather linearly decodable, a learned linear projection over hidden representations can be quite effective. However, this is not always the case, as other literature argues not all the features are linearly decodable \cite{engels2024not}. This makes ReFT potentially less effective where the relevant subspace cannot be manipulated through linear interventions, varying its efficiency from task to task.
	
	\subsection{Layer-wise Analysis}
	Much work has shown how layers can be specialized for certain tasks depending on their depth.
	
	The IEEEtran class file is used to format your paper and style the text. All margins, 
	column widths, line spaces, and text fonts are prescribed; please do not 
	alter them. You may note peculiarities. For example, the head margin
	measures proportionately more than is customary. This measurement 
	and others are deliberate, using specifications that anticipate your paper 
	as one part of the entire proceedings, and not as an independent document. 
	Please do not revise any of the current designations.
	
	\section{Prepare Your Paper Before Styling}
	Before you begin to format your paper, first write and save the content as a 
	separate text file. Complete all content and organizational editing before 
	formatting. Please note sections \ref{AA}--\ref{SCM} below for more information on 
	proofreading, spelling and grammar.
	
	Keep your text and graphic files separate until after the text has been 
	formatted and styled. Do not number text heads---{\LaTeX} will do that 
	for you.
	
	\subsection{Abbreviations and Acronyms}\label{AA}
	Define abbreviations and acronyms the first time they are used in the text, 
	even after they have been defined in the abstract. Abbreviations such as 
	IEEE, SI, MKS, CGS, ac, dc, and rms do not have to be defined. Do not use 
	abbreviations in the title or heads unless they are unavoidable.
	
	\subsection{Units}
	\begin{itemize}
		\item Use either SI (MKS) or CGS as primary units. (SI units are encouraged.) English units may be used as secondary units (in parentheses). An exception would be the use of English units as identifiers in trade, such as ``3.5-inch disk drive''.
		\item Avoid combining SI and CGS units, such as current in amperes and magnetic field in oersteds. This often leads to confusion because equations do not balance dimensionally. If you must use mixed units, clearly state the units for each quantity that you use in an equation.
		\item Do not mix complete spellings and abbreviations of units: ``Wb/m\textsuperscript{2}'' or ``webers per square meter'', not ``webers/m\textsuperscript{2}''. Spell out units when they appear in text: ``. . . a few henries'', not ``. . . a few H''.
		\item Use a zero before decimal points: ``0.25'', not ``.25''. Use ``cm\textsuperscript{3}'', not ``cc''.)
	\end{itemize}
	
	\subsection{Equations}
	Number equations consecutively. To make your 
	equations more compact, you may use the solidus (~/~), the exp function, or 
	appropriate exponents. Italicize Roman symbols for quantities and variables, 
	but not Greek symbols. Use a long dash rather than a hyphen for a minus 
	sign. Punctuate equations with commas or periods when they are part of a 
	sentence, as in:
	\begin{equation}
		a+b=\gamma\label{eq}
	\end{equation}
	
	Be sure that the 
	symbols in your equation have been defined before or immediately following 
	the equation. Use ``\eqref{eq}'', not ``Eq.~\eqref{eq}'' or ``equation \eqref{eq}'', except at 
	the beginning of a sentence: ``Equation \eqref{eq} is . . .''
	
	\subsection{\LaTeX-Specific Advice}
	
	Please use ``soft'' (e.g., \verb|\eqref{Eq}|) cross references instead
	of ``hard'' references (e.g., \verb|(1)|). That will make it possible
	to combine sections, add equations, or change the order of figures or
	citations without having to go through the file line by line.
	
	Please don't use the \verb|{eqnarray}| equation environment. Use
	\verb|{align}| or \verb|{IEEEeqnarray}| instead. The \verb|{eqnarray}|
	environment leaves unsightly spaces around relation symbols.
	
	Please note that the \verb|{subequations}| environment in {\LaTeX}
	will increment the main equation counter even when there are no
	equation numbers displayed. If you forget that, you might write an
	article in which the equation numbers skip from (17) to (20), causing
	the copy editors to wonder if you've discovered a new method of
	counting.
	
	{\BibTeX} does not work by magic. It doesn't get the bibliographic
	data from thin air but from .bib files. If you use {\BibTeX} to produce a
	bibliography you must send the .bib files. 
	
	{\LaTeX} can't read your mind. If you assign the same label to a
	subsubsection and a table, you might find that Table I has been cross
	referenced as Table IV-B3. 
	
	{\LaTeX} does not have precognitive abilities. If you put a
	\verb|\label| command before the command that updates the counter it's
	supposed to be using, the label will pick up the last counter to be
	cross referenced instead. In particular, a \verb|\label| command
	should not go before the caption of a figure or a table.
	
	Do not use \verb|\nonumber| inside the \verb|{array}| environment. It
	will not stop equation numbers inside \verb|{array}| (there won't be
	any anyway) and it might stop a wanted equation number in the
	surrounding equation.
	
	\subsection{Some Common Mistakes}\label{SCM}
	\begin{itemize}
		\item The word ``data'' is plural, not singular.
		\item The subscript for the permeability of vacuum $\mu_{0}$, and other common scientific constants, is zero with subscript formatting, not a lowercase letter ``o''.
		\item In American English, commas, semicolons, periods, question and exclamation marks are located within quotation marks only when a complete thought or name is cited, such as a title or full quotation. When quotation marks are used, instead of a bold or italic typeface, to highlight a word or phrase, punctuation should appear outside of the quotation marks. A parenthetical phrase or statement at the end of a sentence is punctuated outside of the closing parenthesis (like this). (A parenthetical sentence is punctuated within the parentheses.)
		\item A graph within a graph is an ``inset'', not an ``insert''. The word alternatively is preferred to the word ``alternately'' (unless you really mean something that alternates).
		\item Do not use the word ``essentially'' to mean ``approximately'' or ``effectively''.
		\item In your paper title, if the words ``that uses'' can accurately replace the word ``using'', capitalize the ``u''; if not, keep using lower-cased.
		\item Be aware of the different meanings of the homophones ``affect'' and ``effect'', ``complement'' and ``compliment'', ``discreet'' and ``discrete'', ``principal'' and ``principle''.
		\item Do not confuse ``imply'' and ``infer''.
		\item The prefix ``non'' is not a word; it should be joined to the word it modifies, usually without a hyphen.
		\item There is no period after the ``et'' in the Latin abbreviation ``et al.''.
		\item The abbreviation ``i.e.'' means ``that is'', and the abbreviation ``e.g.'' means ``for example''.
	\end{itemize}
	An excellent style manual for science writers is \cite{b7}.
	
	\subsection{Authors and Affiliations}
	\textbf{The class file is designed for, but not limited to, six authors.} A 
	minimum of one author is required for all conference articles. Author names 
	should be listed starting from left to right and then moving down to the 
	next line. This is the author sequence that will be used in future citations 
	and by indexing services. Names should not be listed in columns nor group by 
	affiliation. Please keep your affiliations as succinct as possible (for 
	example, do not differentiate among departments of the same organization).
	
	\subsection{Identify the Headings}
	Headings, or heads, are organizational devices that guide the reader through 
	your paper. There are two types: component heads and text heads.
	
	Component heads identify the different components of your paper and are not 
	topically subordinate to each other. Examples include Acknowledgments and 
	References and, for these, the correct style to use is ``Heading 5''. Use 
	``figure caption'' for your Figure captions, and ``table head'' for your 
	table title. Run-in heads, such as ``Abstract'', will require you to apply a 
	style (in this case, italic) in addition to the style provided by the drop 
	down menu to differentiate the head from the text.
	
	Text heads organize the topics on a relational, hierarchical basis. For 
	example, the paper title is the primary text head because all subsequent 
	material relates and elaborates on this one topic. If there are two or more 
	sub-topics, the next level head (uppercase Roman numerals) should be used 
	and, conversely, if there are not at least two sub-topics, then no subheads 
	should be introduced.
	
	\subsection{Figures and Tables}
	\paragraph{Positioning Figures and Tables} Place figures and tables at the top and 
	bottom of columns. Avoid placing them in the middle of columns. Large 
	figures and tables may span across both columns. Figure captions should be 
	below the figures; table heads should appear above the tables. Insert 
	figures and tables after they are cited in the text. Use the abbreviation 
	``Fig.~\ref{fig}'', even at the beginning of a sentence.
	
	\begin{table}[htbp]
		\caption{Table Type Styles}
		\begin{center}
			\begin{tabular}{|c|c|c|c|}
				\hline
				\textbf{Table}&\multicolumn{3}{|c|}{\textbf{Table Column Head}} \\
				\cline{2-4} 
				\textbf{Head} & \textbf{\textit{Table column subhead}}& \textbf{\textit{Subhead}}& \textbf{\textit{Subhead}} \\
				\hline
				copy& More table copy$^{\mathrm{a}}$& &  \\
				\hline
				\multicolumn{4}{l}{$^{\mathrm{a}}$Sample of a Table footnote.}
			\end{tabular}
			\label{tab1}
		\end{center}
	\end{table}
	
	Figure Labels: Use 8 point Times New Roman for Figure labels. Use words 
	rather than symbols or abbreviations when writing Figure axis labels to 
	avoid confusing the reader. As an example, write the quantity 
	``Magnetization'', or ``Magnetization, M'', not just ``M''. If including 
	units in the label, present them within parentheses. Do not label axes only 
	with units. In the example, write ``Magnetization (A/m)'' or ``Magnetization 
	\{A[m(1)]\}'', not just ``A/m''. Do not label axes with a ratio of 
	quantities and units. For example, write ``Temperature (K)'', not 
	``Temperature/K''.
	
	\section*{Acknowledgment}
	
	The preferred spelling of the word ``acknowledgment'' in America is without 
	an ``e'' after the ``g''. Avoid the stilted expression ``one of us (R. B. 
	G.) thanks $\ldots$''. Instead, try ``R. B. G. thanks$\ldots$''. Put sponsor 
	acknowledgments in the unnumbered footnote on the first page.
	
	\section{References}
	
	\begin{thebibliography}{00}
		
		\bibitem{vaswani2017attention}
		A. Vaswani, N. Shazeer, N. Parmar, J. Uszkoreit, L. Jones, A. N. Gomez, L. Kaiser, and I. Polosukhin, ``Attention is all you need,'' in \textit{Proc. Advances in Neural Information Processing Systems (NeurIPS)}, 2017, pp. 6000--6010.
		
		\bibitem{brown2020language}
		T. B. Brown, B. Mann, N. Ryder, M. Subbiah, J. D. Kaplan, P. Dhariwal, A. Neelakantan, P. Shyam, G. Sastry, A. Askell, S. Agarwal, A. Herbert-Voss, G. Krueger, T. Henighan, R. Child, A. Ramesh, D. Ziegler, J. Wu, C. Winter, C. Hesse, M. Chen, E. Sigler, M. Litwin, S. Gray, B. Chess, J. Clark, C. Berner, S. McCandlish, A. Radford, I. Sutskever, and D. Amodei, ``Language models are few-shot learners,'' in \textit{Proc. Advances in Neural Information Processing Systems (NeurIPS)}, 2020.
		
		\bibitem{hu2022lora}
		E. J. Hu, Y. Shen, P. Wallis, Z. Allen-Zhu, Y. Li, S. Wang, L. Wang, and W. Chen, ``LoRA: Low-rank adaptation of large language models,'' in \textit{Proc. Int. Conf. on Learning Representations (ICLR)}, 2022.
		
		\bibitem{houlsby2019parameter}
		N. Houlsby, A. Giurgiu, S. Jastrzebski, B. Morrone, Q. de Laroussilhe, A. Gesmundo, M. Attariyan, and S. Gelly, ``Parameter-efficient transfer learning for NLP,'' in \textit{Proc. Int. Conf. on Machine Learning (ICML)}, 2019, pp. 2790--2799.
		
		\bibitem{li2021prefix}
		X. L. Li and P. Liang, ``Prefix-tuning: Optimizing continuous prompts for generation,'' in \textit{Proc. Annual Meeting of the Association for Computational Linguistics (ACL)}, 2021.
		
		\bibitem{wu2024reft}
		Z. Wu, A. Arora, Z. Wang, A. Geiger, D. Jurafsky, C. D. Manning, and C. Potts, ``ReFT: Representation Finetuning for Language Models,'' \textit{arXiv preprint arXiv:2404.03592}, 2024.
		
		\bibitem{geiger2023causal}
		A. Geiger, Z. Wu, C. Potts, T. Icard, and N. D. Goodman, ``Finding Alignments Between Interpretable Causal Variables and Distributed Neural Representations,'' \textit{arXiv preprint arXiv:2303.02536}, 2023.
		
		\bibitem{gurnee2023space}
		W. Gurnee and M. Tegmark, ``Language Models Represent Space and Time,'' \textit{arXiv preprint arXiv:2310.02207}, 2023.
		
		\bibitem{engels2024not}
		J. Engels, E. J. Michaud, I. Liao, W. Gurnee, and M. Tegmark, ``Not All Language Model Features Are Linear,'' \textit{arXiv preprint arXiv:2405.14860}, 2024.
	\end{thebibliography}
	
\end{document}
